\documentclass[aspectratio=169,11pt]{beamer}

\usepackage[brazilian]{babel}
\usepackage[utf8]{inputenc}
\usepackage[T1]{fontenc}
\usepackage{amsmath,amssymb,mathtools}
\usepackage{tikz}
\usetikzlibrary{calc,arrows.meta}

\usetheme{Madrid}
\usecolortheme{default}
\setbeamertemplate{navigation symbols}{}

% ---------- macros ----------
\newcommand{\E}[1]{\mathbb{E}\!\left[#1\right]}
\newcommand{\ProbE}[1]{\Pr\!\left[#1\right]}
\newcommand{\Ind}[1]{\mathbf{1}_{#1}}
\newcommand{\paren}[1]{\left(#1\right)}
\newcommand{\define}[1]{\textbf{#1}}
\newcommand{\tri}[3]{%
  \draw (#1,#2) -- ++(-.45,-.9) -- ++(.9,0) -- cycle;
  \node at (#1,#2-.65) {#3};}

\newtheorem{proposicao}[theorem]{Proposição}
\newtheorem{observacao}[theorem]{Observação}
\theoremstyle{definition}

\title{\textit{Treaps}}
\subtitle{Árvores de busca aleatorizadas}
\author{Jair}
\institute{UFABC}
\date{}

\begin{document}

\begin{frame}
  \titlepage
\end{frame}

% =====================================================
\begin{frame}{O problema}
  \textbf{Dicionário}: manter $S\subset U$ sob
  \begin{itemize}
  \item \textbf{busca}($x$): $x\in S$?
  \item \textbf{inserção}($x$)
  \item \textbf{remoção}($x$)
  \end{itemize}
  \pause
  \vspace{1ex}
  Árvores binárias de busca: custo proporcional à \emph{altura}.
  \begin{itemize}
  \item Inserções em ordem crescente $\Rightarrow$ árvore degenerada,
    altura $n-1$.
  \end{itemize}
  \pause
  \vspace{1ex}
  \textbf{Ideia}: usar aleatoriedade para que a árvore se comporte como
  uma árvore de busca \emph{construída em ordem aleatória}, sem
  depender da ordem das operações.
\end{frame}

% =====================================================
\begin{frame}{Definição}
  Uma \define{treap} é uma árvore binária em que cada nó $v$ tem
  $\mathrm{chave}(v)$ e $\mathrm{prioridade}(v)$ tais que
  \begin{itemize}
  \item é \textbf{árvore binária de busca} nas chaves:
    chaves da subárvore esquerda $<$ $\mathrm{chave}(v)$ $<$ chaves da
    subárvore direita;
  \item é \textbf{heap de máximo} nas prioridades:
    $\mathrm{prioridade}(v)>$ prioridade de todo descendente próprio.
  \end{itemize}
  \vspace{1ex}
  Supomos chaves distintas e prioridades distintas.

  \medskip
  \small
  ``Heap'' no sentido de árvore ordenada por prioridades (não é uma
  heap implícita em vetor: a forma não é determinada por completude).
\end{frame}

% =====================================================
\begin{frame}{Exemplo}
  \centering
  \begin{tikzpicture}[scale=.85,
    no/.style={circle, draw, minimum size=1.0cm, inner sep=0pt},
    rot/.style={font=\itshape\small, inner sep=0pt},
    val/.style={font=\small, inner sep=0pt}]
    \node[no] (j) at (0,0)       {};
    \node[no] (i) at (-3,-1.7)   {};
    \node[no] (m) at (3,-1.7)    {};
    \node[no] (a) at (-4.5,-3.4) {};
    \node[no] (b) at (-1.5,-3.4) {};
    \node[no] (l) at (1.5,-3.4)  {};
    \draw (j) -- (i);
    \draw (j) -- (m);
    \draw (i) -- (a);
    \draw (i) -- (b);
    \draw (m) -- (l);
    \foreach \n/\letra/\valor in {j/j/31, i/b/26, m/m/27, a/a/13, b/i/25, l/\ell/14}{
      \draw (\n.45) -- (\n.-135);
      \node[rot] at ($(\n)+(-0.2,0.2)$) {$\letra$};
      \node[val] at ($(\n)+(0.2,-0.2)$) {$\valor$};
    }
  \end{tikzpicture}

  \small letras = chaves \quad números = prioridades \quad
  em-ordem: $a,b,i,j,\ell,m$
\end{frame}

% =====================================================
\begin{frame}{Unicidade}
  \begin{proposicao}
    Dado um conjunto de pares (chave, prioridade) com chaves distintas
    e prioridades distintas, existe uma \textbf{única} treap que o
    representa.
  \end{proposicao}
  \pause
  \textbf{Ideia} (indução no número de nós): a \emph{heap} obriga a
  raiz a ser o nó de maior prioridade; a busca obriga a subárvore
  esquerda a ter exatamente as chaves menores que a da raiz e a direita
  as maiores; cada subárvore é uma treap, única por indução.
  \pause
  \medskip

  \textbf{Consequências}
  \begin{itemize}
  \item A forma depende só de chaves e prioridades, \emph{não} da
    ordem de inserção.
  \item Equivale a inserir os nós em ordem \emph{decrescente} de
    prioridade numa ABB vazia.
  \end{itemize}
\end{frame}

% =====================================================
\begin{frame}{Rotações}
  \centering
  \begin{tikzpicture}[scale=.8, >={Stealth[length=3mm,width=3.5mm]},
      no/.style={draw,rectangle,minimum size=5mm,inner sep=1pt,fill=white}]
    \node[no] (y)  at (0,0)  {$y$};
    \node[no] (x)  at (-1,-1) {$x$};
    \draw (y) -- (x);
    \draw (y) -- (1,-1);   \tri{1}{-1}{$C$}
    \draw (x) -- (-1.6,-2); \tri{-1.6}{-2}{$A$}
    \draw (x) -- (-0.4,-2); \tri{-0.4}{-2}{$B$}
    \draw[->] (2.3,-0.8) to[bend left=40]
      node[above] {direita em $y$} (4.2,-0.8);
    \draw[->] (4.2,-1.3) to[bend left=40]
      node[below] {esquerda em $x$} (2.3,-1.3);
    \begin{scope}[xshift=6.5cm]
      \node[no] (x2) at (0,0)  {$x$};
      \node[no] (y2) at (1,-1) {$y$};
      \draw (x2) -- (y2);
      \draw (x2) -- (-0.6,-1); \tri{-0.6}{-1}{$A$}
      \draw (y2) -- (0.4,-2);  \tri{0.4}{-2}{$B$}
      \draw (y2) -- (1.8,-2);  \tri{1.8}{-2}{$C$}
    \end{scope}
  \end{tikzpicture}

  \medskip
  \begin{itemize}
  \item A ordem em-ordem $A,x,B,y,C$ é preservada.
  \item Só a subárvore $B$ muda de pai: custo $O(1)$.
  \end{itemize}
\end{frame}

% =====================================================
\begin{frame}{Operações}
  \begin{description}
  \item[Busca] como em ABB; custo $\propto$ altura.
  \pause
  \item[Inserção de $x$]
    \begin{enumerate}
    \item busca por $x$; se falhar, insere $x$ como folha com uma
      prioridade \textbf{sorteada};
    \item enquanto $x$ tem prioridade maior que a do pai $y$: rotação
      que sobe $x$.
    \end{enumerate}
    A \emph{heap} só pode ser violada entre $x$ e seu pai, e cada
    rotação preserva a busca. Rotações $\le$ profundidade de $x$.
  \pause
  \item[Remoção de $x$] enquanto $x$ não é folha, rotaciona $x$ com o
    filho de \textbf{maior} prioridade (que passa a ser o pai de $x$
    e, portanto, vence o outro filho); depois remove a folha.
    Rotações $\le$ altura da subárvore de $x$.
  \end{description}
\end{frame}

% =====================================================
\begin{frame}{Aleatoriedade}
  \begin{itemize}
  \item Prioridade de cada nó: sorteada uniformemente em $[0,1]$ na
    inserção, independente das demais (reinserir uma chave sorteia
    uma prioridade \emph{nova}).
  \item Com probabilidade $1$, as prioridades são distintas.
  \end{itemize}
  \pause
  \begin{alertblock}{Adversário oblívio}
    A sequência de operações é escolhida \emph{independentemente} das
    prioridades.

    Se o adversário visse as prioridades, poderia remover sempre os nós
    de menor prioridade, e as restantes deixariam de ser i.i.d.\
    uniformes.
  \end{alertblock}
  \pause
  Pela unicidade, após qualquer sequência oblívia de operações a
  árvore é \textbf{a treap do conjunto atual, com prioridades i.i.d.\
    uniformes}.

  $\Rightarrow$ basta analisar uma treap com $n$ nós e prioridades
  i.i.d.\ uniformes.
\end{frame}

% =====================================================
\begin{frame}{Profundidade de um nó}
  Seja $x_k$ o $k$-ésimo menor elemento de $S$.
  \begin{itemize}
  \item $x_j\prec x_k$: $x_j\neq x_k$ é \emph{ancestral próprio} de
    $x_k$.
  \item $I_j=\Ind{[x_j\prec x_k]}$.
  \item \textbf{Profundidade}: $D_k=\sum_{j\neq k}I_j$ (em arestas).
  \item \textbf{Altura}: $\max_k D_k$.
  \item $X(j,k)=\{x_j,\dots,x_k\}$ (para $j<k$; $X(k,j)=X(j,k)$).
  \end{itemize}
  \pause
  \begin{proposicao}
    Para $j\neq k$: $x_j$ é ancestral próprio de $x_k$ $\iff$ $x_j$
    tem a \textbf{maior prioridade} em $X(j,k)$.
  \end{proposicao}
\end{frame}

% =====================================================
\begin{frame}{Prova da proposição}
  Seja $a$ o ancestral comum mais baixo de $x_j$ e $x_k$ (todo nó é
  ancestral de si mesmo, só nesta prova).
  \begin{enumerate}
  \item[(i)] $a\in X(j,k)$: se $a\notin\{x_j,x_k\}$, eles estão em
    subárvores distintas de $a$, e a chave de $a$ fica entre as deles.
    \pause
  \item[(ii)] $X(j,k)\subseteq$ subárvore de $a$: acima de $a$, $x_j$
    e $x_k$ estão do mesmo lado de cada ancestral $b$, logo qualquer
    $x_i$ com chave entre as deles também. A busca por $x_i$ segue o
    caminho de $x_j$ até $a$.
    \pause
  \item[(iii)] Pela \emph{heap}, $a$ tem prioridade maior que a de
    todos os outros nós da sua subárvore; portanto é o único de maior
    prioridade em $X(j,k)$.
  \end{enumerate}
  \pause
  Como $j\neq k$: $x_j\prec x_k\iff a=x_j\iff x_j$ tem a maior
  prioridade em $X(j,k)$. \hfill $\square$
\end{frame}

% =====================================================
\begin{frame}{Profundidade esperada}
  Prioridades i.i.d.\ contínuas $\Rightarrow$ todo elemento de $X(j,k)$
  é o de maior prioridade com a mesma probabilidade, e
  $|X(j,k)|=|k-j|+1$:
  \[
    \ProbE{I_j=1}=\frac{1}{|k-j|+1}.
  \]
  \pause
  \[
    \E{D_k}
    =\sum_{j<k}\frac1{k-j+1}+\sum_{j>k}\frac1{j-k+1}
    =H_k+H_{n-k+1}-2 .
  \]
  \pause
  Com $H_m\le1+\ln m$:
  \[
    \E{D_k}\le2H_n-2\le2\ln n=O(\log n)\quad\text{para todo }k.
  \]
  \medskip
  Esperança do número de nós no caminho até $x_k$:
  $H_k+H_{n-k+1}-1$.
\end{frame}

% =====================================================
\begin{frame}{Alta probabilidade: independência}
  Queremos Chernoff, que exige independência.

  \begin{alertblock}{Os $I_j$ \emph{não} são todos independentes}
    $n=3$, $k=2$:
    $\ProbE{x_1\prec x_2,\;x_3\prec x_2}
    =\ProbE{\text{a prioridade de }x_2\text{ é a menor}}=\tfrac13$,
    mas $\ProbE{x_1\prec x_2}\ProbE{x_3\prec x_2}=\tfrac14$.
  \end{alertblock}
  \pause
  \vspace{1ex}
  Separamos $D_k=L_k+R_k$ com
  $L_k=\sum_{j<k}I_j$ e $R_k=\sum_{j>k}I_j$.

  \begin{lemma}
    Fixado $k$, os $I_j$ com $j<k$ são \textbf{independentes}, com
    $\ProbE{I_j=1}=\frac1{k-j+1}$. O mesmo vale para $j>k$.
  \end{lemma}
\end{frame}

% =====================================================
\begin{frame}{Prova do lema}
  \begin{itemize}
  \item $y_i$ = prioridade de $x_{k-i+1}$, $i=1,\dots,k$.
  \item Pela proposição: $I_{k-i+1}=1\iff y_i>\max\{y_1,\dots,y_{i-1}\}$.
  \pause
  \item $\rho_i\in\{1,\dots,i\}$ = posto de $y_i$ entre $y_1,\dots,y_i$.
  \item A aplicação
    $\text{permutação}\mapsto(\rho_1,\dots,\rho_k)$ é injetora
    (reconstrói-se a permutação inserindo $y_1,y_2,\dots$ nas posições
    dadas pelos postos) e os dois conjuntos têm $k!$ elementos:
    é bijetora.
  \pause
  \item Logo $\rho_1,\dots,\rho_k$ são independentes, $\rho_i$ uniforme
    em $\{1,\dots,i\}$.
  \item $I_{k-i+1}=\Ind{[\rho_i=i]}$ depende só de $\rho_i$:
    independentes, com probabilidade $1/i$. \hfill $\square$
  \end{itemize}
\end{frame}

% =====================================================
\begin{frame}{Teorema: altura logarítmica}
  \begin{theorem}
    Numa treap com $n\ge2$ nós e prioridades i.i.d.\ uniformes, com
    probabilidade pelo menos $1-2n^{-7}$ a profundidade de todo nó é
    menor que $2\mathrm e^2\ln n<15\ln n$.
  \end{theorem}
  \pause
  Em particular, a altura é $O(\log n)$ com probabilidade
  $1-O(n^{-7})$, e, nesse instante, cada operação de dicionário tem
  custo $O(\log n)$ com essa probabilidade.

  \medskip
  \small
  ($n\ge2$ garante $\ln n>0$; para $n=1$ a desigualdade estrita
  falharia.)
\end{frame}

% =====================================================
\begin{frame}{Prova do teorema (1/2)}
  Fixe $k$. $L_k$ é soma de $k-1$ Bernoullis independentes de
  parâmetros $1/i$, $2\le i\le k$, e
  \[
    \E{L_k}=H_k-1\le\ln n .
  \]
  \pause
  Para $s>0$, usando $1+x\le\mathrm e^x$:
  \[
    \E{\mathrm e^{sL_k}}
    =\prod_{i=2}^{k}\paren{1+\frac{\mathrm e^s-1}{i}}
    \le\exp\paren{(\mathrm e^s-1)\,\E{L_k}}
    \le\exp\paren{(\mathrm e^s-1)\ln n}.
  \]
  \pause
  Markov em $\mathrm e^{sL_k}$, com $s=2$ e limiar $\mathrm e^2\ln n$:
  \[
    \ProbE{L_k\ge\mathrm e^2\ln n}
    \le\mathrm e^{-2\mathrm e^2\ln n}\,\mathrm e^{(\mathrm e^2-1)\ln n}
    =n^{-(\mathrm e^2+1)}<n^{-8}.
  \]
\end{frame}

% =====================================================
\begin{frame}{Prova do teorema (2/2)}
  O mesmo vale para $R_k$. Se $L_k,R_k<\mathrm e^2\ln n$, então
  $D_k<2\mathrm e^2\ln n$.
  \pause
  \[
    \ProbE{\exists k:\;D_k\ge2\mathrm e^2\ln n}
    \le 2n\cdot n^{-8}=2n^{-7}.
  \]
  \pause
  Fora desse evento, todo nó tem profundidade
  $<2\mathrm e^2\ln n<15\ln n$ (pois $2\mathrm e^2\approx14{,}78$).

  Altura $<15\ln n$ em arestas ($\le15\ln n+1$ em nós). \hfill $\square$
\end{frame}

% =====================================================
\begin{frame}{Sequências de operações}
  O teorema refere-se a uma treap num instante fixo.

  \medskip
  Numa sequência oblívia de $m$ operações, seja $N\ge2$ o maior número
  de nós ao longo da sequência.
  \begin{itemize}
  \item Cada estado é uma treap com prioridades i.i.d.\ uniformes.
  \item Em cada estado $\E{L_k}\le\ln n\le\ln N$: a prova vale com
    $\ln N$ no lugar de $\ln n$.
  \item Subaditividade sobre os $m$ estados:
    \[
      \text{todas as operações custam}<15\ln N+1
      \ \text{com prob.}\ \ge1-2m\,N^{-7}.
    \]
  \item Para $m=O(N^2)$: probabilidade $1-O(N^{-5})$.
  \end{itemize}
\end{frame}

% =====================================================
\begin{frame}{Resumo}
  \begin{itemize}
  \item Treap = ABB nas chaves + heap de máximo nas prioridades;
    \textbf{forma única} dado o conjunto de pares.
  \item Inserção e remoção por \textbf{rotações}, $O(1)$ cada.
  \item Prioridades aleatórias (adversário oblívio) $\Rightarrow$
    árvore como ABB de ordem aleatória.
  \item $x_j\prec x_k\iff x_j$ tem a maior prioridade em $X(j,k)$.
  \item $\E{D_k}\le2\ln n$.
  \item Altura $<15\ln n$ com probabilidade $\ge1-2n^{-7}$, via
    independência de cada lado + MGF/Markov (Chernoff).
  \end{itemize}
\end{frame}

\end{document}
